\section{The recursive carrier theorem}\label{sec:recursive-carriers}

Composite contraction reduces the uniform problem to the geometry of linear
spaces contained in the terminal bad locus.  We determine that locus after
reduction in each geometric fiber and show that all modular descendants merge
into one linear carrier.

For \(r\geq6\), put \(n=r-1\), and let the ground field have
characteristic \(p>0\).  In divided-power coordinates define
\begin{equation}\label{eq:maximal-lucas-carrier}
 \cM^{\max}_{r,p}=
 \PP\!\left\langle e_j:
 0\leq j\leq r-1,\quad
 \binom{r-2}{j}=\binom{r-2}{j-1}=0\pmod p\right\rangle ,
\end{equation}
where \(\binom nk=0\) when \(k\notin\{0,\ldots,n\}\).
An empty span denotes the empty projective scheme.  Let
\(\cP_r\) be the rank-at-most-two consecutive catalecticant scheme in
\(\PP(\Gamma^{r-1}E)\).

\begin{definition}[Reduced terminal bad scheme]\label{def:terminal-bad-scheme}%
\coverage{absent}%
For a cubic pencil with Pl\"ucker coordinates \(z_0,\ldots,z_5\), put
\[
 \Phi_z(t,s)=z_0+z_1(t+s)+z_2(t^2+ts+s^2)+z_3ts
              +z_4ts(t+s)+z_5t^2s^2
\]
and let \(K\subset\mathbb Z[z_0,\ldots,z_5]\) be the ideal
\begin{align*}
(&3z_4^2-9z_2z_5-z_3z_5, z_3z_4-3z_1z_5,
 z_3^2-9z_0z_5,\\
 &z_2z_3-z_1z_4+z_0z_5, z_1z_3-3z_0z_4,
 3z_1^2-9z_0z_2-z_0z_3).
\end{align*}
The polynomial \(\Phi_z\) is the off-diagonal ordered-pair equation of the
pencil: the diagonal factor has already been divided out.  The ideal \(K\)
is the saturated coefficient-elimination ideal of its symmetric residual
factorization; the anti-invariant collision component is removed before
elimination.  This identification is the computer-assisted statement checked
by Certificate SC.

For
\[
 z(c)=(c_2c_4-c_3^2,-c_1c_4+c_2c_3,c_1c_3-c_2^2,
 c_0c_4-c_1c_3,-c_0c_3+c_1c_2,c_0c_2-c_1^2)
\]
and
\[
 D=\det\begin{pmatrix}
 c_0&c_1&c_2\\c_1&c_2&c_3\\c_2&c_3&c_4
 \end{pmatrix},
\]
define in characteristic \(p\)
\[
 I_{\mathrm{res},p}
 =\sqrt{\bigl((z^*K)\otimes\F_p:D^\infty\bigr)},
 \qquad
 \mathcal B_{5,p}^{\mathrm{red}}
 =V(D)\cup V(I_{\mathrm{res},p}).
\]
This is the \emph{reduced terminal bad scheme}.  Thus the scheme being
decomposed is fixed before its component calculation; it is not inferred
from a finite-field census.
\end{definition}

The next two propositions and one theorem prove the unconditional geometric
layer: they find every reduced terminal component, merge all modular descendants, and
select the components that persist up the contraction tower.  On a first
pass, read their statements and the four-row component table in
Proposition~\ref{prop:recursive-component-selection}, then resume at
Theorem~\ref{thm:recursive-carrier}.  The intervening elimination and
characteristic-wise substitutions prove exhaustiveness and add no hypothesis
to the finite-depth theorem.

\begin{proposition}[Computer-assisted terminal decomposition]\label{prop:reduced-terminal-carrier}%
\coverage{absent}%
\evidence{certificate-stable-components}%
The reduced terminal bad scheme of
Definition~\ref{def:terminal-bad-scheme} is the irredundant union of the prime
cubic \(V(D)\) and one residual prime.  Away from characteristics two and
three the residual prime is the projected Veronese surface
\[
 [u:v:w]\longmapsto [6u^2:3uv:v^2+2uw:3vw:6w^2],
 \autoeqtag{eq:terminal-residual-param}
\]
the locus of syndrome quartics that are scalar multiples of squares of binary
quadratics.  The coordinates in \eqref{eq:terminal-residual-param} are the
bottom \(c_i\); writing \(z_i=\binom4ic_i\) for the plain coefficients of the
associated quartic \(f_c=\sum_{i=0}^4z_iX^{4-i}Y^i\), the parametrization gives
\(f_c=6(uX^2+vXY+wY^2)^2\), so the same surface is the
plain-coefficient Veronese in rescaled divided-power coordinates.
In characteristic two it is the plane \(V(c_0,c_4)\).  In
characteristic three it is the cone with \(c_2\) free and prime ideal
\[
 (c_1c_3-c_0c_4,\ c_1^3-c_0^2c_3,\
   c_0c_3^2-c_1^2c_4,\ c_3^3-c_1c_4^2).
 \autoeqtag{eq:terminal-char3-prime}
\]
There is no flat reduced integral model having exactly these residual
fibers.

For a base-point-free separable terminal pencil, lying outside this scheme is
equivalent to geometric \(S_3\) monodromy of its cubic map.  Thus
Proposition~\ref{prop:r5-count} and Lemma~\ref{lem:r5-branch} apply to every
terminal point outside \(\mathcal B_{5,p}^{\mathrm{red}}\).
\end{proposition}

\begin{proof}
The elimination has a concrete geometric source.  On
\(\operatorname{Gr}(2,\Gamma^3E)\), use the Plücker coordinates and
ordered-pair equation of Definition~\ref{def:terminal-bad-scheme}.
For the Hankel matrix
\[
 H_c=\begin{pmatrix}c_0&c_1&c_2&c_3\\c_1&c_2&c_3&c_4\end{pmatrix},
\]
the signed maximal minors give the map explicitly:
\[
 z(c)=(c_2c_4-c_3^2,\,-c_1c_4+c_2c_3,\,c_1c_3-c_2^2,
 c_0c_4-c_1c_3,\,-c_0c_3+c_1c_2,\,c_0c_2-c_1^2).
\]
After the diagonal and collision factors already charged as deletion
divisors are saturated away, the terminal bad locus is the pullback, along
the Hankel Plücker map \(c\mapsto z(c)\), of the factorization locus of
\(\Phi_z\).  The involution \((t,s)\leftrightarrow(s,t)\) leaves three
possibilities.  Symmetric factors give the residual component below;
exchanged factors give \(D=0\); and an anti-invariant factor is proportional
to \(s-t\).  In the last case the only nonzero Plücker coordinates are
\(z_2=-\mu\) and \(z_3=3\mu\), and hence
\[
 z_0z_5-z_1z_4+z_2z_3=(-\mu)(3\mu)=-3\mu^2.
\]
Thus it
survives only in characteristic three and is then the diagonal collision
already removed.  Thus the two displayed primes are the only possible
noncollision components before small-characteristic specialization.

The exact saturated bridge can also be read without the elimination file.
Let \(J=z^*K\) in the notation of
Definition~\ref{def:terminal-bad-scheme}.
If \(V\) denotes the seven-generator residual ideal printed below, exact
reduction gives \(J\subseteq V\), \(3D\,V\subseteq J\), and
\((J:D^\infty)=V\) over \(\mathbb Q\).  The factor \(3\) is necessary over
\(\mathbb Z\), which is precisely why the characteristic-three fiber is
checked separately.

Outside characteristics two and three, elimination of the ordered residual
roots gives the prime kernel of \eqref{eq:terminal-residual-param}; one
convenient generating set is
\begin{align*}
&2c_3^3-3c_2c_3c_4+c_1c_4^2,\\
&6c_2c_3^2-9c_2^2c_4+2c_1c_3c_4+c_0c_4^2,\\
&2c_1c_3^2-3c_1c_2c_4+c_0c_3c_4,\quad
  c_0c_3^2-c_1^2c_4,
\end{align*}
together with their images under \(c_i\leftrightarrow c_{4-i}\).  Direct
reduction gives the plane in characteristic two and
\eqref{eq:terminal-char3-prime} in characteristic
three.  The ordered-Hessian calculation sends its complementary ruling to
rank one.

Here exhaustiveness is an exact elimination statement, not a visual
factorization heuristic.  Certificate SC starts from the terminal
ordered-root incidence ideal, saturates the collision factors, eliminates
the root variables, and verifies the two generic primes and their
characteristic-\(2,3\) fibers by Gröbner reduction.  Its public Python and
Singular inputs are in
\texttt{supplement/evidence/stable-components/}; every generator and
remainder used for this conclusion is hash-pinned there.

Each displayed residual ideal is prime:
\eqref{eq:terminal-residual-param} is a projective parametrization with prime
kernel, the characteristic-two ideal is linear, and
\eqref{eq:terminal-char3-prime} is the
toric prime of the monomial quartic after projection from \(e_2\).  None is
contained in \(V(D)\): use \((3,0,1,0,3)\), which is the image of
\((u,v,w)=(1,0,1)\) and has \(D=8\), together with \(e_2\) and
\((1,0,1,0,0)\), respectively.  Thus the intersection with \((D)\) is an
irredundant reduced decomposition and has no embedded prime.  The residual
degree is four generically and one in characteristic two.  Constancy of the
Hilbert polynomial in a flat projective family rules out the final assertion.
The exact eliminations and the characteristic-\(2,3,5\) reductions are
replayed by the stable-component certificates described in
Section~\ref{sec:verification}.

Finally, a base-point-free separable cubic map has transitive geometric
monodromy \(C_3\) or \(S_3\).  Its off-diagonal ordered-pair curve is
reducible exactly in the cyclic case, and the noncollision reducibility is
precisely the saturated residual factorization encoded by \(K\).  A pencil
outside the displayed scheme therefore has geometric monodromy \(S_3\),
which proves the last clause.
\end{proof}

\begin{definition}[Coherent polar flag]
A coherent polar flag from redundancy \(r\) to redundancy five is a sequence
\[
 f=f_r,\qquad f_{j-1}=\iota_{\lambda_j}f_j
 \quad (j=r,\ldots,6),
\]
with distinct geometric markers \(\lambda_j\), such that every contraction is
nonzero and every earlier marker remains a forbidden root for the lower
witness system.  A coherent lift reverses this sequence using
Equation~\eqref{eq:lift}.
\end{definition}

\begin{proposition}[Maximal Lucas union]\label{prop:maximal-lucas-union}%
\coverage{absent}%
\imports{wang-wu-hu}%
Every normal-rational-curve nucleus encountered in a coherent polar flag,
and every later lift of it, is contained in
\(\cM^{\max}_{r,p}\).  Their reduced union equals that single linear
scheme.  A polar line not contained in the union meets it in degree at most
one.
\end{proposition}

\begin{proof}
For the order-\(a\) nucleus of the degree-\(d\) curve, its coordinate support
is contained in
\[
 \widehat S_{d,p}=\{i:\binom di=0\pmod p\},
 \autoeqtag{eq:lucas-zero-support}
\]
because the nucleus conditions include the row \(d\).  If a support is
\(S\subset\{0,\ldots,d\}\), one coherent lift has support
\[
 L_d(S)=\{j:(j=d+1\text{ or }j\in S),\
                 (j=0\text{ or }j-1\in S)\}.
 \autoeqtag{eq:lucas-support-lift}
\]
This operation preserves inclusions.  Pascal's identity gives
\(L_d(\widehat S_{d,p})\subset\widehat S_{d+1,p}\).
Induction therefore contains every earlier descendant in the last one.
Conversely, the top-order nucleus at \(d=r-2\) occurs in the union, and its
one-step lift is exactly \eqref{eq:maximal-lucas-carrier}.  The union is
therefore one linear scheme.  Pulling its linear ideal back to a projective
line gives nonzero linear forms unless the line is contained, proving the
degree bound.
\end{proof}

The endpoint arithmetic in Proposition~\ref{prop:maximal-lucas-union} has a
known projective-subline form.  Wang--Wu--Hu
\cite[Prop.~11]{WangWuHu2026} prove that the Lucas block with
\(k=p^a+1\) is nonempty over \(\F_{p^e}\) exactly when \(a\mid e\), and
then its blocks are projective sublines.  We use this criterion at that
exact boundary.  The adjacent-zero carrier
\eqref{eq:maximal-lucas-carrier}, its nesting through coherent lifts, and
its split-squarefree Hankel incidence are separate questions addressed here.

\begin{definition}[Reduced recursively contained carrier]
Take the reduced union of the irreducible upper components whose generic
coherent contraction row space is contained in a component of
\(\mathcal B_{5,p}^{\mathrm{red}}\), together with every coherent lift of an
intermediate normal-rational-curve nucleus.  This is the \emph{reduced
recursively contained carrier}.  Fixed-gcd strata and isolated marker
collisions are excluded: they are handled by lower packages and deletion
divisors, respectively.
\end{definition}

\begin{theorem}[Recursive carrier]
\label{prop:recursive-component-selection}%
\coverage{fragment}%
The reduced recursively contained carrier generated by the terminal carrier
and all intermediate nuclei is exactly
\begin{equation}\label{eq:recursive-carrier}
                       \cP_r\cup\cM^{\max}_{r,p}.
\end{equation}
Fixed linear gcds retain their quadratic-graph package, and marker collisions
remain deletion divisors; neither is included in
\eqref{eq:recursive-carrier}.
\end{theorem}

\begin{proof}
Retain \(m\) markers.  If \(h\) is the coefficient row of their product,
the bottom syndrome is
\[
 h\operatorname{Cat}_{m,4}(f),\qquad
 \operatorname{Cat}_{m,4}(f)=
 \begin{pmatrix}
 a_0&a_1&a_2&a_3&a_4\\
 a_1&a_2&a_3&a_4&a_5\\
 \vdots&\vdots&\vdots&\vdots&\vdots\\
 a_m&a_{m+1}&a_{m+2}&a_{m+3}&a_{m+4}
 \end{pmatrix}.
 \autoeqtag{eq:marker-catalecticant}
\]
Over an algebraic closure, products of distinct linear forms are dense in
the projective space of binary forms.  Deleting the projective kernel of the
linear map does not change its image closure.  Hence the attainable bottom
syndromes close to the projectivized row space of
\eqref{eq:marker-catalecticant}.  This row space is
irreducible.  If it lies in the finite reduced union of
Proposition~\ref{prop:reduced-terminal-carrier}, its generic point lies in
one component, so the whole row space lies in that component.  This is a
geometric-marker statement; no density of the finite set of rational marker
tuples is used.  The density and finite-component-selection implication are
also kernel checked by the declarations listed in the formalization ledger.

The remaining component check is summarized by
\[
\begin{array}{c|c|c}
\text{component}&\text{possible linear part}&\text{upper lift}\\ \hline
V(D)&\text{arbitrary}&\cP_r\\
\text{Veronese }(p\ne2,3)&\text{point}&\text{rank one}\\
\text{wild cone }(p=3)&\text{ruling through }e_2&\text{rank one}\\
V(c_0,c_4)\ (p=2)&\text{binary-plane subspace}&
   \cM^{\max}_{r,2}
\end{array}
\]
Here is the calculation behind the four rows.  Substitute
\(c_i=xd_i+yd_{i+1}\).  For \(V(D)\), the coefficients in \(x,y\) are
exactly the \(3\times3\) minors of
\[
 \begin{pmatrix}d_0&d_1&d_2&d_3\\d_1&d_2&d_3&d_4\\
 d_2&d_3&d_4&d_5\end{pmatrix},
\]
the consecutive Fano identity.  For the generic residual prime their unique
minimal prime is the ideal of the \(2\times2\) minors of
\(\bigl(\begin{smallmatrix}d_0&\cdots&d_4\\d_1&\cdots&d_5\end{smallmatrix}\bigr)\),
so the upper syndrome has rank one; geometrically this is also the fact that
the projected Veronese contains no line.  The same coefficient ideal in
characteristic three is again this rank-one ideal, verifying directly that
every positive-dimensional ruling of the wild cone through \(e_2\) lifts
into \(\cP_r\).  In characteristic two it has exactly the two minimal primes
\[
 (d_0,d_1,d_4,d_5)
 \quad\text{and}\quad
 (d_3d_4+d_2d_5,d_1d_4+d_0d_5,d_3^2+d_1d_5,
 d_2d_3+d_0d_5,d_2^2+d_0d_4,d_1d_2+d_0d_3).
\]
The second lies in the persistent \(3\times3\)-minor ideal; the first is the
binary-plane descendant.  Successive overlaps cut it to the binary line,
then the central point, then the empty projective space, and
\eqref{eq:lucas-support-lift} identifies precisely these spaces as Lucas
lifts.  On a nucleus, \eqref{eq:lucas-support-lift} and
Proposition~\ref{prop:maximal-lucas-union} then place every surviving lift in
the maximal modular carrier.

The coherent-Fano part of Certificate SC verifies these consecutive-row
substitutions, including the characteristic-three ruling and the
characteristic-two successive supports, against the same incidence ideal.
This exhausts every component and proves \eqref{eq:recursive-carrier} by
induction on the flag length.
\end{proof}

\begin{theorem}[Arbitrary-redundancy carrier theorem]
\label{thm:recursive-carrier}%
\coverage{absent}%
\uses{thm:simultaneous-marker-escape}%
\evidence{certificate-stable-components}%
\proves{recursive-carrier-classification}%
Let \(r\geq6\), and let \(q\) be a prime power with
\begin{equation}\label{eq:uniform-threshold}
 q\geq Q_r^*:=6r-16+\left\lfloor2\sqrt{6r-18}\right\rfloor.
\end{equation}
Every split-free redundancy-\(r\) syndrome lies in
\(\cP_r\cup\cM^{\max}_{r,p}\).  If \(p>r-1\), the modular term is empty,
and the projective deep-hole syndromes are exactly the persistent tangent
and conjugate-secant families.  Their cardinality is
\[
                         \frac{q(q+1)^2}{2},
\]
and their orbit law is \(T/T^{r-1}\) modulo inversion and coefficient
Frobenius, with tangent cocycle \(z\mapsto z+(r-1)u\).  In characteristic
two the carrier containment holds under
\(q\geq6r-22+\lfloor2\sqrt{6r-24}\rfloor\).
\end{theorem}

\begin{proof}
Proposition~\ref{prop:recursive-component-selection} identifies the complete
reduced carrier.  Theorem~\ref{thm:simultaneous-marker-escape}, with (s=0),
gives a split squarefree kernel member at every point outside it.  Its strict
genus-one inequality is exactly \eqref{eq:uniform-threshold}; no intermediate
lower package is assumed.

If \(p>r-1\), no interior entry of Pascal row \(r-2\) vanishes, so
\(\cM^{\max}_{r,p}=\varnothing\).  Rank one and rational secants are shallow;
the remaining rank-two points are precisely the tangent and conjugate-secant
families and are split-free.  The threshold lies in the
Seroussi--Roth--Dür radius range, so split-free and deep coincide.  The count
and orbit law are the persistent calculation already used at R6 and R7.
\end{proof}

The carrier theorem is fiberwise and reduced.  It makes no assertion about nilpotent
structure in a chosen integral incidence model, and in small characteristic
it classifies the only possible carrier of additional split-free points, not
the arithmetic of every point on that carrier.  The simultaneous-marker
theorem proves escape from its complement directly.

\begin{proposition}[Upper-level radius gate]\label{prop:upper-radius}%
\coverage{fragment}%
\imports{duer,seroussi-roth}%
For the ranges used below,
\[
\begin{aligned}
 \rho(\PRS(q-7))&=7 &&(q\geq43),\\
 \rho(\PRS(q-8))&=8 &&(q\geq53),\\
 \rho(\PRS(q-9))&=9 &&(q\geq59).
\end{aligned}
\]
More generally, the radius is \(r-1\) under the full-support,
large-characteristic hypotheses of Theorem~\ref{thm:main}.
\end{proposition}

\begin{proof}
The dual GRS code has dimension \(r\).  In each displayed range, and also
under \eqref{eq:uniform-threshold}, one has
\[
                         r\leq\left\lfloor q/2\right\rfloor+2.
\]
Seroussi--Roth \cite[Thm.~1]{SeroussiRoth1986} therefore gives
nonextendability of the corresponding normal rational curve; its stated
even-field dimension-three exception is irrelevant because \(r\geq8\).
Dür's completeness--covering-radius equivalence \cite{Dur1994}, in the
form recorded by Kaipa \cite[Sec.~IV and Prop.~4(1)]{Kaipa2017}, gives
the asserted radius.  Thus the promotion from split-free syndromes to deep
holes in R8--R10 uses one explicit imported gate.
\end{proof}
